The rapid advancement of diffusion-based image generation models such as Midjourney and Stable Diffusion has significantly increased the realism of synthetic images, raising concerns regarding misinformation, identity misuse, and digital security. This paper presents a comparative analysis of five deepfake detection models: CNNDetect, EfficientNet, CLIP with LoRA adaptation, DIRE, and FIRE. The study evaluates both zero-shot generalization and fine-tuned performance using a unified subset derived from OpenFake, containing 8,989 in-distribution images: 3,000 real LAION images, 2,989 Midjourney-generated images, and 3,000 Stable Diffusion-generated images. To support fair comparison, all models were evaluated on identical dataset splits, while fine-tuning strategies were adapted to each model's architecture and computational requirements. Additionally, the models are evaluated on an out-of-distribution test set containing 450 real images and 900 fake images from generators unseen during training. Experimental results reveal significant differences in generalization capability, adaptability to diffusion-generated images, and robustness to distribution shift.
